\section*{ID6301: Relation: φ ≈ 1.618034}
\paragraph{Metadata.}
PiID: 9601212285993 | RTEID: RTE:P24468.1823:T2.6169 | UUID: 86e04565-d6c2-5368-a260-9dcbfe17cb23

\paragraph{Canonical Statement.}
Done — I created a 3D visualization that treats the d’Alembertian as a pyramid (time as the apex, spatial directions as the base) and overlaid a phase winding around the base whose winding number equals the golden ratio \textbackslash{}(\textbackslash{}varphi\textbackslash{}!\textbackslash{}approx\textbackslash{}!1.618034\textbackslash{}).

\paragraph{Canonical Equation.}
\[
\varphi\!\approx\!1.618034
\]

\paragraph{Dictionary.}
Done — I created a 3D visualization that treats the d’Alembertian as a pyramid (time as the apex, spatial directions as the base) and overlaid a phase winding around the base whose winding number equals the golden ratio \textbackslash{}( \textbackslash{}!

\paragraph{Encyclopedia.}
Relation: φ ≈ 1.618034 is a scaling / order-of-magnitude relation, written φ ≈ 1.618034. It expresses a scaling or proportionality, capturing how the quantity grows with its parameters. It is built from φ. Within the Trawin Zero Point registry it serves as one element of the framework's mathematical narrative.
